\documentclass[17pt]{beamer}

\usepackage[utf8]{inputenc}
\usepackage[turkish]{babel}
\usepackage[T1]{fontenc}

\usepackage{amsmath}
\usepackage{amssymb}
\usepackage{amsthm}
\usepackage{enumerate}
\usepackage{tikz}
\usepackage{transparent}
\usepackage{xcolor}

\usetheme{Warsaw}
\usecolortheme{default}
\definecolor{nigdeyesili_acik}{RGB}{3, 150, 166}
\definecolor{nigdeyesili_koyu}{RGB}{143, 209, 217}
\setbeamercolor{structure}{fg=nigdeyesili_acik}

\author[Dr. Mehmet CANEVİ]{Arş.~Gör.~Dr.~M.~Canevi\inst{1}}

\institute{
    \inst{1}%
    Bilgisayar Mühendisliği\\
    Mühendislik Fakültesi
}
    
\date[2025] {Ders Notları, Ocak 2025}
    
\logo{\transparent{0.4}\includegraphics[width = 20mm]{logo}}
    
\title[Ders 8] {Dosyalar ile Çalışma}
\begin{document}
%%%%%%%%%%%%%%%%%%%%%%%%%%%%%%%%%%%%%%%%%%%%%%%%%%%%%%%%%%%%%%%%%%%%%%%%%%%%%%%%
\frame{\titlepage}
\begin{frame}[fragile]{İçidekiler}
    \tableofcontents
\end{frame}
%%%%%%%%%%%%%%%%%%%%%%%%%%%%%%%%%%%%%%%%%%%%%%%%%%%%%%%%%%%%%%%%%%%%%%%%%%%%%%%%
\section{Dosyalar}
\begin{frame}[fragile]{Dosya Okuma}
    \begin{lstlisting}
        FILE* dosya=fopen("dosya_adi.txt", "r");
        char metin[100];
        while(fgets(metin, 100, dosya)) {
            printf("%s", metin);
        }
        fclose(dosya);\end{lstlisting}
\end{frame}
%%%%%%%%%%%%%%%%%%%%%%%%%%%%%%%%%%%%%%%%%%%%%%%%%%%%%%%%%%%%%%%%%%%%%%%%%%%%%%%%
\begin{frame}[fragile]{Dosyaya Yazma}
    \begin{lstlisting}
        FILE* dosya=fopen("dosya_adi.txt", "w");
        fprintf(dosya, "deneme");
        fclose(dosya);\end{lstlisting}
\end{frame}
%%%%%%%%%%%%%%%%%%%%%%%%%%%%%%%%%%%%%%%%%%%%%%%%%%%%%%%%%%%%%%%%%%%%%%%%%%%%%%%%
\section{Uygulama}
\begin{frame}[fragile]{Soru}
    \begin{alertblock}{Soru-1}
        Bir sinyali dosyaya yazan kodu yazınız.
    \end{alertblock}
    \begin{alertblock}{Soru-2}
        Birden çok sinyali dosyaya yazan kodu yazınız.
    \end{alertblock}
    \begin{alertblock}{Soru-3}
        CSV dosyasını okuyan kodu yazınız.
    \end{alertblock}
    \begin{alertblock}{Soru-4}
        Adres defteri okuyan kodu yazınız.
    \end{alertblock}
\end{frame}
%%%%%%%%%%%%%%%%%%%%%%%%%%%%%%%%%%%%%%%%%%%%%%%%%%%%%%%%%%%%%%%%%%%%%%%%%%%%%%%%
\end{document}